
\setlength{\parindent}{0em}

\begin{abstract}

Vehicular Ad Hoc Networks (VANETs) are an important
component of Intelligent Transportation Systems (ITS),
enabling wireless communication among vehicles and
roadside infrastructure to support road safety,
traffic management, and transportation efficiency.
This paper presents a literature-based technical
review of VANET architectures, communication models,
routing protocols, simulation frameworks, and
emerging Vehicle-to-Everything (V2X) technologies.
Vehicle-to-Vehicle (V2V), Vehicle-to-Infrastructure
(V2I), and other V2X communication modes are examined
alongside proactive, reactive, and geographic routing
approaches. A comparative analysis of published
studies highlights important trade-offs associated
with vehicle mobility, network connectivity,
communication overhead, and routing reliability.
The review also discusses simulation frameworks,
including SUMO, OMNeT++, Veins, and NS-3, and
performance metrics such as packet delivery ratio,
end-to-end delay, and throughput. Modern communication
technologies, including IEEE 802.11p, LTE-V2X,
IEEE 802.11bd, and 5G NR-V2X, are reviewed.
Major research challenges involving security,
privacy, channel congestion, scalability, and
interoperability are identified, together with
potential future directions. As a literature-based
study, this paper does not report original
simulation experiments or newly measured results.
Instead, it synthesizes existing research to
provide a structured understanding of VANET
technologies, limitations, and research opportunities.

\end{abstract}
